\section{Introduction}  \label{sec:intro}

This paper is about the complexity class $\MIP^*$ of multiprover interactive proof systems with
entangled quantum provers---the quantum version of the classical class
$\MIP$. Classically, the study of $\MIP$ has had far-reaching
implications in theoretical computer science. In complexity theory, the proof by
Babai, Fortnow, and Lund~\cite{BFL91} that $\MIP = \NEXP$ was the
direct antecedent of the PCP theorem~\cite{ALM+98,AS98}, a seminal
result which is the foundation of the modern theory of hardness of approximation. In cryptography, the $\MIP$ model was introduced
to allow for information-theoretic zero-knowledge proofs~\cite{BGKW88}, and more recently $\MIP$ protocols have become essential building blocks
in designing delegated computation schemes (see e.g.~\cite{KRR14}). These implications alone
would be a sufficient motivation for considering the quantum class
$\MIP^*$, but remarkably, the study of $\MIP^*$ is also deeply related to
long-standing questions in the foundations of quantum
mechanics regarding the nature of quantum entanglement. Indeed, the
$\MIP^*$ model itself was anticipated by the \emph{nonlocal games} or
\emph{Bell tests} introduced in the
work of John Bell~\cite{Bel64}, who was in turn inspired by the thought
experiment proposed by Einstein, Podolsky, and Rosen~\cite{EPR35}.
These nonlocal games have had applications to quantum
cryptography~\cite{Eke91, MY98, Col06}, delegated quantum
computation~\cite{RUV13}, and more.

Even though the class $\MIP$ is now well-understood, it has proven difficult to
determine the computational power of $\MIP^*$. A priori, it is not
even clear that  $\MIP^*$ contains $\MIP$, since adding entanglement could increase or decrease the power of
the proof system. This is because the added resource of entanglement
can make it easier for dishonest provers to cheat the verifier. Indeed, Cleve et al.~\cite{CHTW04} showed that for proof systems
based on so-called XOR games (where the verifier's decision can only
depend on the XOR of the provers' answer bits), the quantum class
$\oplus\MIP^* \subseteq \EXP$, whereas classically $\oplus\MIP =
\NEXP$. In particular, this result implied that the classical
$\oplus\MIP$ protocol for $\NEXP$ of H{\aa}stad~\cite{Has97} could not be sound against entangled
provers.
In spite of this, Ito and Vidick~\cite{IV12,Vid16} were able to show that $\NEXP \subseteq
\MIP^*$, by proving that a different classical protocol \emph{is}
sound against entanglement. Note that the protocol of~\cite{Vid16} is
\emph{identical} to a protocol shown to be unsound by Cleve et al.,
except in that it uses 3 provers rather than 2 (the protocol is played
by choosing a random subset of 2 provers from the 3). This illustrates
the subtleties of dealing with entangled provers.

With the lower bound $\NEXP \subseteq \MIP^*$ established, a natural
follow-up question is whether $\MIP^*$ is \emph{strictly} more powerful
than $\MIP$. Indeed, it was long known that some $\MIP^*$ protocols
possess a uniquely quantum property called \emph{self-testing}, which
has no direct analog in the classical setting. Roughly speaking, an $\MIP^*$ protocol is
a self-test for a particular entangled state $\ket{\psi}$ if only
provers using states close to $\ket{\psi}$ can achieve close to optimal
success in the protocol. In such a protocol, observing that the
provers succeed with nearly optimal probability \emph{certifies} that
they share a state close to the target state $\ket{\psi}$. The germ of
this idea came from the work of Bell~\cite{Bel64}, who studied the types of
bipartite correlations that could be obtained from measuring an
entangled state called the EPR state, which had been introduced by
Einstein, Podolsky, and Rosen~\cite{EPR35}. Bell gave a protocol where
provers using the EPR state could succeed with a greater probability
than purely classical provers, and subsequent works of Tsirelson~\cite{Tsi80}, and Summers and
Werner~\cite{SW88} showed that (a variant of) Bell's protocol certifies
the EPR state in the sense of self-testing.


In order to prove stronger lower bonds on $\MIP^*$, the post-Ito-Vidick phase of
$\MIP^*$ research aimed to use this self-testing property to design protocols
for problems in Hamiltonian complexity, the quantum analog of the
theory of $\NP$-completeness. In Hamiltonian complexity, the complexity class
$\QMA$ plays the role of $\NP$; it is the set of problems for which
there exists a quantum witness state that can be efficiently checked
by a polynomial-time quantum verifier. Problems in $\QMA$ seemed like a natural match for
the powers of $\MIP^*$ as one could potentially construct a protocol for $\QMA$ by
designing a self-test for accepting witness states of some $\QMA$-complete
problem. The connection between $\MIP^*$ and $\QMA$ was also well
motivated from the point of view of the ``quantum PCP'' research
program, which strives to find quantum analogues of the classical PCP
theorem. In the classical setting, the PCP theorem can be viewed as a scaled-down
version of $\MIP^* = \NEXP$, showing that there exists an $\MIP^*$
protocol for 3SAT (and thus for all of $\NP$) with $O(\log(n))$-sized
messages. Drawing inspiration from this, Fitzsimons and Vidick~\cite{FV15}
stated a ``quantum games PCP conjecture'': that there should exist an
$\MIP^*$ protocol with $\log(n)$-sized messages for the local
Hamiltonian problem, and thus for the class $\QMA$. This
was proved by Natarajan and Vidick~\cite{NV18a} in 2018 with a 7-prover protocol. Along the way to achieving
this goal,~\cite{NV18a} developed a highly efficient self-test for
high-dimensional entangled states: their ``quantum low-degree test''
is a self-test for $n$ EPR pairs with only $O(\log(n))$ communication.

Already, the result of~\cite{NV18a} is strong evidence that $\MIP^*
\neq \MIP$, since it is believed that $\QMA \neq \NP$. But, at the same
time, several other works showed that even larger separations were
possible in the regime of subconstant soundness gaps. Here there are
results in two settings. For $\MIP^*$ with a soundness gap scaling
inverse-exponentially (i.e. $1/\exp(n)$) in the instance size,
Ji~\cite{Ji17} showed a protocol for $\NEEXP$: nondeterministic
\emph{doubly}-exponential time, and a subsequent work by Fitzsimons,
Ji, Vidick, and Yuen~\cite{FJVY19} showed protocols for
non-deterministic iterated exponential time (e.g. $\NTIME(2^{2^n})$)
with a correspondingly small soundness gap (e.g. $2^{-C \cdot
  2^n}$). In the ``gapless'' case, Slofstra~\cite{Slo16,Slo19} showed that given
a description of an $\MIP^*$ protocol, determining whether there exists an entangled
strategy that succeeds with probability exactly $1$ is undecidable by any Turing machine.

These results hint at the full power of $\MIP^*$ but are not
conclusive, as it is not unusual for quantum complexity classes to
increase significantly in power when a numerical precision parameter
is allowed to shrink. For instance, $\QIP$ (quantum interactive proofs
with a single prover) with an exponentially small
gap is equal to $\EXP$~\cite{IKW12}, while $\QIP$ with a polynomial gap is equal to
$\IP = \PSPACE$. Likewise, $\QMA$ with exponentially small
gap (known as $\mathsf{PreciseQMA}$) is known to be equal to
$\PSPACE$~\cite{FL18}, while $\QMA$ is contained $\PP$, and $\QMA(k)$
($\QMA$ with multiple unentangled Merlins) with exponentially small gap is
equal to $\NEXP$~\cite{Per12}, whereas in the constant-gap
regime the best known lower bound is that
$\QMA(k) \supseteq \QMA$. Moreover, even the $\QMA$ lower bound for
$\MIP^*_{\log}$ obtained by~\cite{NV18b} holds for $7$ provers only;
with $2$ provers, the best known lower bound for $\MIP^*_{\log}$ is
$\NP = \MIP_{\log}$~\cite{NV18a}. Could it be that 2-prover $\MIP^*$ is
equal to $\MIP$, with entanglement providing no advantage at all?

This paper conclusively answers this question in the negative. Our main result (\Cref{thm:main-informal}) is to show that $\MIP^*$ contains $\NEEXP$, with only
two provers and with a constant completeness-soundness gap. This is
establishes the first known unconditional separation between $\MIP^*$
and $\MIP$ in the constant-gap regime: previously, such a separation
was known only assuming $\QMA \neq \NP$, and only in the scaled-down
setting of logarithmic-sized messages.

\begin{theorem}[\Cref{thm:main} in the body]
  \label{thm:main-informal}
  There is a two-prover, one-round $\MIP^*$ protocol for the $\NEEXP$-complete
  problem $\succinctsquared$ with completeness $1$, soundness $1/2$, and
  question and answer length $\poly(n)$.
\end{theorem}

As a corollary of \Cref{thm:main-informal}, we obtain a lower bound on the hardness
of approximation for the entangled value $\omega^*$ of a nonlocal game.
\begin{corollary}\label{cor:main-games-informal}
  There exists a constant $c < 1$ such that given a two-prover
  nonlocal game $\game$ of size $N$, the problem of deciding
  whether $\omega^*(\game) = 1$ or $\omega^*(\game) \leq 1/2$,
  promised one of the two holds, is $\NTIME(2^{N^{\log^{-c} N}})$-hard.
\end{corollary}
For two-player games, the best prior lower bound was
$\NP$~\cite{NV18b}. The lower bound achieved in
\Cref{cor:main-games-informal} is
stronger as for any $c < 1$, the function $2^{N^{\log^{-c} N}}$ is
superpolynomial.

\paragraph{Techniques.}

Our construction, inspired by~\cite{Ji17}
and~\cite{FJVY19}, involves ``compression'': we show how to take an $\MIP$ protocol for $\NEEXP$
with exponentially-long questions and answers (the ``big'' protocol),
and simulate it by an $\MIP^*$ protocol with polynomial-sized messages (the ``small''
protocol). However, the techniques we use to achieve our compression
are quite different. We
eschew the Hamiltonian-complexity ideas that were used in previous
works, and in particular the use of history states. In our protocol,
honest provers need only share a quantum resource state of (exponentially
many) EPR pairs, together with a \emph{classical} assignment to the
$\NEEXP$ instance being tested. The use of history states was the main
barrier preventing previous works from applying to the case of
two provers. 

We
divide compression into two steps: \emph{question compression} and
\emph{answer compression}. We
achieve question compression by a technique which we call
\emph{introspection}, in which we command the provers to perform
measurements on their shared EPR pairs whose outcomes are pairs of questions
from the ``big'' protocol. To force the provers to sample their
questions honestly, we use a variant of the quantum low-degree test
from~\cite{NV18a}, which certifies Pauli measurements on exponentially
many EPR pairs using messages of only polynomial size. A crucial
challenge is to prevent each prover from learning the other prover's
sampled question, since this would destroy the soundness of the
``big'' protocol. To achieve this, we use the \emph{``data-hiding''} properties of quantum measurements
in incompatible bases: if a set of qubits is measured in the Pauli $X$-basis,
this ``erases'' all information about $Z$-basis
measurements. This means that if Alice samples her question by
measuring her half of a block of EPR pairs in the $Z$-basis, then her question can
be hidden from Bob by forcing him (via self-testing) to measure his
half of the EPR pairs in the $X$-basis. Interestingly, our data-hiding scheme
does \emph{not} operate in a black-box way on the ``big'' protocol, but
rather makes essential use of its structure. In particular, we start
with a ``big'' protocol based on a scaled-up version of a PCP
construction using the low-degree test, where the question
distribution consists of pairs of random points in a vector space and
affine subspaces containing them. The linear structure of the vector
space is essential for our data-hiding procedure to work.

Our approach to answer compression is more standard, essentially using composition with a
classical PCP of proximity. Here, the verifier asks the provers
to compute a PCP proof that their ``big'' answers satisfy the success
conditions of the protocol, and verifies this PCP proof by reading an
exponentially smaller number of bits. Care is needed to deal
with entanglement between the provers. The first, fundamental
challenge we face is that the success condition of the ``big''
protocol is a function of \emph{both} provers' answers. Thus, to
compute a PCP proof that the condition is satisfied, one of the
provers must have access to both provers' answers. Classically, this
is achieved using the technique of oracularization, in which one prover
receives \emph{both} provers' questions and is checked for consistency
against the other prover, which only receives a single question. In
the entangled setting, this oracularization procedure is sound, but not
necessarily complete. This is because oracularization requires that each prover, if given the
\emph{other} prover's question, could predict its answer with
certainty, even though this answer is obtained from a nondeterministic
quantum measurement. In our protocol, we are able to use
oracularization because honest provers always use
a maximally entangled state, which they measure with projective
measurements that pairwise commute for every pair of questions asked
in the game. While this commutation requirement is restrictive, it
still permits non-trivial quantum behavior; indeed, the linear system games used by Slofstra~\cite{Slo19}
involve similar commutation conditions.


The second challenge is to ensure that the PCP of proximity we use for
composition is itself sound against entanglement. We achieve this by
performing a further step of composition: we ask the provers to encode
their PCP proof in the low-degree code and verify it with the
low-degree test, which is known to be sound even against entangled
provers~\cite{NV18b}. This technique was introduced in the $\QMA$
protocol of~\cite{NV18a} in
order to perform energy measurements on the provers' state.



\paragraph{Implications and future work}
We believe that our work opens up several exciting directions for
further progress. For the complexity theorist, the most obvious future
direction is to obtain even stronger lower bounds on $\MIP^*$ by
iterating our protocol, as in~\cite{FJVY19}. At the most basic level,
we could imagine taking our $\MIP^*$ protocol for $\NEEXP$ and
performing a further layer of question compression and answer compression
on it, thus obtaining an $\MIP^*$ protocol with
logarithmic message size for $\NEEXP$, or, scaling up, an $\MIP^*$
protocol with polynomial message size for
$\ntime(2^{2^{2^{\poly(n)}}})$. By further iterating question
reduction and answer reduction $k$ times, we could obtain potentially
obtain lower bounds of
$\ntime(\underbrace{2^{\cdot^{\cdot^{n}}}}_{k})$ on $\MIP^*$ while
retaining a constant completeness-soundness gap. The main obstacle to
achieving such results is that the question compression procedure
developed in this paper is tailored to a special distribution of
questions (that of the $\MIP_{\exp}$ protocol for $\NEEXP$), whereas our
answer compression procedure produces protocols whose question
distribution is not of this form.

Assuming that this obstacle can be surmounted, we could aspire to a
more ambitious goal: a general ``gap-preserving compression procedure'' for some
subclass of $\MIP^*$ protocols, which we may label ``compressible'' protocols. Such a procedure would consist of a
Turing  machine that takes as input any compressible $\MIP^*$
protocol $\game$, and generates a new compressible protocol $\game'$ with exponentially
smaller message size, but approximately the same entangled value. It
was shown by~\cite{FJVY19} that the existence of such a compression
procedure for the set of \emph{all} $\MIP^*$ protocols would imply
that $\MIP^*$ contains the set of all computable languages,
and moreover that there exists an undecidable language in
$\MIP^*$. These consequences would continue to hold as long as the set
of compressible protocols contains a family of protocols solving
problems in $\NTIME(f(n))$, where $f(n)$ is a growing function of
$n$.

Showing that $\MIP^*$ contains undecidable languages would be
significant not just for complexity theory but also for the
foundations of quantum mechanics, as it would resolve a long-standing
open problem known as \emph{Tsirelson's problem}. Tsirelson's problem
asks whether two notions of quantum nonlocality are equivalent: the
\emph{tensor-product model}, in which two parties Alice and Bob each
act on their respective factor of a tensor-product Hilbert space
$\calH_{\reg{Alice}} \ot \calH_{\reg{Bob}}$, and the
\emph{commuting-operator model}, in which both parties act on a common
Hilbert space $\calH$, but the algebra of Alice's measurement
operators must commute with Bob's, and vice versa. It was shown by
Slofstra~\cite{Slo16} that in the ``zero-error'' setting, these two
models differ: there are quantum correlations which can be
\emph{exactly} achieved in the commuting-operator model but not in the
tensor product model. Surprisingly, showing that $\MIP^*$ contains
undecidable languages would imply that the two models are separated
even in the bounded-error setting: it would imply that there exist
correlations that can be achieved in the commuting-operator model that cannot
even be approximated (up to constant precision) in the tensor-product
model. The reason for this implication is that if the two models are
indistinguishable up to bounded error, then there exists a Turing machine
that can decide any language in $\MIP^*$ and is guaranteed to halt. 
This observation, which is folklore in the community, follows from the completeness of the non-commutative
sum of squares hierarchy for the commuting-operator model, as
documented in~\cite{FJVY19}. Showing a separation between the two
models would have significant mathematical consequences as well, as it
would yield a negative answer to the long-standing Connes' embedding problem.

In addition to these connections to complexity and mathematical
physics, we hope that our results will have applications in other
areas such as to delegated computation or quantum cryptography. In particular,
our use of introspection is reminiscent of ideas used in quantum
randomness expansion, where randomness generated by measuring EPR
pairs is used to generate questions for a nonlocal game. Could our
results improve on the infinite randomness expansion protocol of
Coudron and Yuen~\cite{CY14}?
\paragraph{Acknowledgements}
We thank Henry Yuen for many useful conversations about the idea of
``introspecting'' interactive proof protocols, which inspired us to
start this project. AN is also grateful to the Simons Institute for
the hospitable environment of the Summer Cluster on Challenges in
Quantum Computation during which these conversations where held. We thank Thomas Vidick for his guidance and advice.  We thank Ryan O'Donnell and
Ryan Williams for a succinct review of the literature on the
complexity of succinct (succinct) 3Sat and $\mathsf{NE(E)XP}$. We are
also grateful to Zhengfeng Ji for several useful discussions, especially regarding the
consequences of recursively composing our protocol with itself.

AN was partially supported by NSF grant CCF-1452616.  JW was partially
supported by ARO contract W911NF-17-1-0433. Both authors acknowledge funding provided by the Institute for Quantum Information and Matter, an NSF Physics Frontiers Center (NSF Grant PHY-1733907).


%%% Local Variables:
%%% mode: latex
%%% TeX-master: "main"
%%% End:
%%%---------- close: introduction

%%%%%%%%%%% jump to overview
%%%---------- open: overview
%!TEX root = main.tex

